\documentclass[11pt]{article}
\usepackage[margin=1in]{geometry}
\usepackage{amsmath,amssymb,amsthm,enumitem,hyperref}
\hypersetup{colorlinks=true,urlcolor=blue,linkcolor=blue}
\newtheorem{theorem}{Theorem}
\newtheorem{lemma}{Lemma}
\title{The norm--supernorm relation for the size ensemble}
\author{Principia Math In Progress: MathDB 359712}
\date{21 September 2026}
\begin{document}
\maketitle
\begin{abstract}
We prove that the product of the reciprocal norm and reciprocal supernorm
partition sums of size $n$ tends to one. The logarithmic derivative of the
supernorm generating function gives a positive recurrence. Two successive
bounds on its coefficients turn the known cumulative asymptotic into the
required individual asymptotic.
\end{abstract}

\section*{Task 1: the original problem}
For a partition $\lambda$ of size $|\lambda|=n$, write
\[
 N(\lambda)=\prod_i\lambda_i,\qquad
 \widehat N(\lambda)=\prod_i p_{\lambda_i},\qquad
 v_n=\sum_{|\lambda|=n}\frac1{N(\lambda)},\qquad
 u_n=\sum_{|\lambda|=n}\frac1{\widehat N(\lambda)},
\]
where $p_j$ is the $j$th prime, $p_1=2$, and $u_0=v_0=1$.
The intended $W_{\rm size}$ in the source is the reciprocal norm sum $v_n$.

\begin{theorem}\label{thm:main}
As $n\to\infty$,
\[
                     u_n\sim\frac{e^\gamma}{n},
 \qquad\text{and consequently}\qquad v_nu_n\longrightarrow1.
\]
\end{theorem}

The product assertion is Conjecture 3.8 in both versions of \cite{LS}.
Its equivalent size-supernorm assertion is Conjecture 3.5 of the preprint
and Conjecture 3.7 of the journal final. They are counted as one target.
The separate perimeter conjecture is not used in this proof. We use three
established inputs from \cite{LS}: the cumulative theorem, the prime number
theorem, and Lehmer's reciprocal norm theorem:
\begin{equation}\label{eq:inputs}
 C_n:=\sum_{j=0}^n u_j\sim e^\gamma\log n,
 \qquad p_n\sim n\log n,
 \qquad v_n\sim e^{-\gamma}n.
\end{equation}
See respectively Theorem 3.4, Lemma 4.4, and Theorem 1.1 there.

\section*{The logarithmic-derivative kernel}
As an identity of formal power series,
\[
 F(z)=\sum_{n\ge0}u_nz^n
     =\prod_{d\ge1}(1-z^d/p_d)^{-1}.
\]
Every coefficient depends on finitely many factors. Formal logarithmic
differentiation therefore gives
\begin{equation}\label{eq:kernel}
 \frac{zF'(z)}{F(z)}=\sum_{k\ge1}a_kz^k,
 \qquad a_k=\sum_{d\mid k}\frac{d}{p_d^{k/d}}.
\end{equation}
Equating coefficients yields the exact positive recurrence
\begin{equation}\label{eq:rec}
                    nu_n=\sum_{k=1}^n a_k u_{n-k}.
\end{equation}

\begin{lemma}\label{lem:kernel}
As $k\to\infty$,
\[
 a_k=\frac{k}{p_k}+O(k^{-1})\sim\frac1{\log k}.
\]
Consequently $a_k$ is bounded and
$\sum_{k\le x}a_k=O(x/\log x)$ as $x\to\infty$.
\end{lemma}
\begin{proof}
Write $R_k=a_k-k/p_k\ge0$ and put $r=k/d\ge2$ in its sum.
For $2\le r\le\sqrt k$, the elementary bound $p_d\ge d+1>d$ gives
\[
 \frac{d}{p_d^r}\le\left(\frac{r}{k}\right)^{r-1}.
\]
For $k\ge4$, the $r=2$ term is at most $2/k$. The terms with
$3\le r\le\sqrt k$, after dropping the divisibility restriction, total at most
\[
 \sum_{r=3}^\infty k^{-(r-1)/2}
   =\frac{1/k}{1-1/\sqrt k}\le\frac2k.
\]
For $r>\sqrt k$, the corresponding $d$ are less than $\sqrt k$.
There are at most $\sqrt k$ such integers, each term is at most
$\sqrt k\,2^{-\sqrt k}$, and thus
\begin{equation}\label{eq:R}
 0\le R_k\le\frac4k+k\,2^{-\sqrt k}=O(k^{-1}).
\end{equation}
The prime number theorem now proves the kernel asymptotic and boundedness.
To bound its partial sums, the $k\le\sqrt x$ terms total $O(\sqrt x)$;
the remaining terms are each $O(1/\log x)$, giving $O(x/\log x)$.
\end{proof}

\section*{Two bounds and the limiting constant}
Boundedness of $a_k$ and \eqref{eq:inputs} first imply from \eqref{eq:rec}
\[
 nu_n=O(C_{n-1})=O(\log n),\qquad u_n=O(\log n/n).
\]
Split \eqref{eq:rec} at $k=n/2$. For $k\le n/2$, this initial bound gives
$u_{n-k}=O(\log n/n)$ uniformly, so Lemma~\ref{lem:kernel} implies
\[
 \sum_{k\le n/2}a_ku_{n-k}
 =O\left(\frac{\log n}{n}\right)
  O\left(\frac n{\log n}\right)=O(1).
\]
For $k>n/2$, the kernel bound is uniformly $O(1/\log n)$, whence
\[
 \sum_{n/2<k\le n}a_ku_{n-k}
 =O\left(\frac1{\log n}\right)C_{\lceil n/2\rceil-1}=O(1).
\]
Consequently
\begin{equation}\label{eq:improved}
                              u_n=O(1/n).
\end{equation}
Use \eqref{eq:improved} in the first half of the same recurrence to improve
its contribution to
\begin{equation}\label{eq:small}
 \sum_{k\le n/2}a_ku_{n-k}
 =O\left(\frac1n\right)O\left(\frac n{\log n}\right)=o(1).
\end{equation}
In the second half, $a_k\log n\to1$ uniformly for $n/2<k\le n$:
indeed $a_k\log k\to1$ and $\log k/\log n\to1$ uniformly there.
Since every $u_j\ge0$, this gives
\[
 \sum_{n/2<k\le n}a_ku_{n-k}
 =\frac{1+o(1)}{\log n}C_{\lceil n/2\rceil-1}
 \longrightarrow e^\gamma.
\]
Together with \eqref{eq:small}, recurrence \eqref{eq:rec} proves
$nu_n\to e^\gamma$. Multiplication by Lehmer's asymptotic in
\eqref{eq:inputs} proves $v_nu_n\to1$.\hfill$\square$

\section*{Scope and verification}
All positive integer parts and their multiplicities are unrestricted.
The proof uses the already proved cumulative theorem without differentiating
its error term. Formal coefficient operations require no boundary convergence.
The companion verifier independently enumerates partitions, compares a
product expansion with the exact recurrence, checks the kernel remainder
bound with rational arithmetic, and supplies numerical examples.
Finite checks are not a proof of the limit. Review is by the authoring
agent only, without independent peer review or a proof assistant.

\section*{Proposed follow-up Tasks 2--5}
\begin{enumerate}[label=Task \arabic*:,start=2,leftmargin=4.6em]
\item Obtain an explicit remainder estimate for $nu_n-e^\gamma$, with a
stated constant and threshold.
\item Prove a coefficient transfer theorem for positive recurrences of the
form \eqref{eq:rec} when $a_k\sim\alpha/\log k$ and
$\sum_{j\le n}u_j\sim c(\log n)^\alpha$, for $\alpha,c>0$.
\item Determine the leading constant for the reciprocal supernorm sum over
partitions of size $n$ with distinct parts.
\item For each fixed $\beta>1$, determine the leading asymptotic of
$\sum_{|\lambda|=n}\widehat N(\lambda)^{-\beta}$, including its constant.
\end{enumerate}

\begin{thebibliography}{1}
\bibitem{LS} J. C. Lagarias and C. Sun,
\emph{Asymptotics of reciprocal supernorm partition statistics},
The Ramanujan Journal \textbf{65} (2024), 487--520.
\href{https://doi.org/10.1007/s11139-024-00893-8}{Journal version (DOI)}.
Preprint: \href{https://arxiv.org/abs/2308.15303}{arXiv:2308.15303v1}.
\end{thebibliography}
\end{document}
